The studies performed by Amir et. al~\cite{Amir2014,Amir2014_2} provided movies of microscopic
images which showcased 3 states in total:
An initial state, a bent state and the final relaxed state.
To effectively estimate the positions of the bacterium at these time points, we manually modified
the provided images by amplifying light shadings of the green cellular envelope.
This is displayed in Figure~\ref{fig:supplement-amir-progression} (A-C) as the bright green cellular
shape.
Our position extraction algorithm provides the discretized vertices of the bacterium (shown in red).
Figures (D-F) show renders of the corresponding simulation snapshots where dark gray symbolizes the
area in which the confinement of the tube is employed as described in
Section~\ref{sec:application-1-amir}.
The system is specified by the estimated parameters and a range of constants which can be obtained
from the experimental conditions.
These are listed in Table~\ref{tab:supplement-amir-constants}.
The number of vertices $N_\text{vertices}$ should be chosen large enough such that the bending of
the rod can be properly captured.
Values which are too small will inadvertently lead to undesired interactions with the confinement
condition imposed by the narrow tube.\\
To fully model this system, a force exerted by the flowing fluid is taken into account.
The flow profile was derived from a standard laminar flow along a tube with Neumann boundary
conditions~\cite{Wolschin2021}.
This profile was experimentally measured and confirmed~\cite{Amir2014}.
Furthermore, the publication explicitly stated values for the size of the flow domain north of the
confined tube $L_{y,\text{domain}}$ which leaves us only with the strength of the flow force as the
remaining estimated parameter.

\begin{figure}[H]
    \centering
    \begin{tikzonimage}[width=0.3\textwidth]
        {figures/crm_amir/progressions/step4-000010.png}
        \node at (0.02,17pt)[anchor=north west, white]{\textbf{A}};
    \end{tikzonimage}%
    \hspace{0.01\textwidth}%
    \begin{tikzonimage}[width=0.3\textwidth]
        {figures/crm_amir/progressions/step4-000017.png}
        \node at (0.02,17pt)[anchor=north west, white]{\textbf{B}};
    \end{tikzonimage}%
    \hspace{0.01\textwidth}%
    \begin{tikzonimage}[width=0.3\textwidth]
        {figures/crm_amir/progressions/step4-000034.png}
        \node at (0.02,17pt)[anchor=north west, white]{\textbf{C}};
    \end{tikzonimage}\\
    \vspace{0.01\textwidth}%
    \begin{tikzonimage}[width=0.3\textwidth]
        {figures/crm_amir/result-full/0000000000.png}
        \node at (0.02, 0.98)[anchor=north west]{\textbf{D}};
    \end{tikzonimage}%
    \hspace{0.01\textwidth}%
    \begin{tikzonimage}[width=0.3\textwidth]
        {figures/crm_amir/result-full/0000000001.png}
        \node at (0.02, 0.98)[anchor=north west]{\textbf{E}};
    \end{tikzonimage}%
    \hspace{0.01\textwidth}%
    \begin{tikzonimage}[width=0.3\textwidth]
        {figures/crm_amir/result-full/0000000002.png}
        \node at (0.02, 0.98)[anchor=north west]{\textbf{F}};
    \end{tikzonimage}%
    \caption{
        (A-C) Images by Amir et. al~\cite{Amir2014_2} which have been manually segmented.
        The red line within the mask represents the extracted positions.
        (D-F) Snapshots of the numerical simulation showing the rod in its initial, fully bent and
        fully relaxed state.
    }
    \label{fig:supplement-amir-progression}
\end{figure}

\begin{table}[H]
    \begin{tabularx}{\textwidth}{l r@{} @{\hspace{0.3em}}l X}
        \textbf{Constant} & \multicolumn{2}{l}{\textbf{Value}} & \textbf{Description}\\
        \toprule
        $t_\text{bend}$ & $7$&$\unit{\second}$ & Interval with active flow\\
        $t_\text{relax}$ & $17$&$\unit{\second}$ & Relaxation interval after deactivating flow\\
        $h$ & $10.2$&$\unit{\pixel\per\micro\metre}$ & Conversion factor between pixels and lengths\\
        $N_\text{vertices}$ & $20$ && Number of vertices use for the discretized rod\\
        $\Delta x$ & $0.8$&$\unit{\micro\metre}$ & Displacement uncertainty used to calculate the profile-likelihood\\
        $L_\text{y,domain}$ & $59.22$ & $\unit{\micro\metre}$ & Vertical size of the domain
        extending further than the shown images.\\
            & $604$ & $\unit{\pixel}$\\
        $L_\text{tube}$ & $19.61$ & $\unit{\micro\metre}$ & Size of the tube which confines the bacterium.\\
            & $200$ & $\unit{\pixel}$\\
        \bottomrule
    \end{tabularx}
    \caption{Constants used for the numerical simulation.}
    \label{tab:supplement-amir-constants}
\end{table}

%###################################################################################################
% \paragraph*{Analysis \& Usage of heterogeneous Parameter Value (TODO)}
% \label{sec:supplement-reuse-estimated-parameters}
% In order to reuse this model, we also provide the option to assign growth rates by sampling from a
% normal distribution with user-defined parameters.
